\documentclass[11pt]{article}
\usepackage[margin=1in]{geometry}
\usepackage{amsmath}
\usepackage{amssymb}
\usepackage{booktabs}
\usepackage{hyperref}
\usepackage{listings}
\usepackage{xcolor}
\usepackage{fancyhdr}

\lstset{
    basicstyle=\ttfamily\small,
    backgroundcolor=\color{gray!10},
    frame=single,
    breaklines=true
}

\pagestyle{fancy}
\fancyhf{}
\rhead{June 5, 2026}
\lhead{Cross-Correlation Analysis Plan}
\cfoot{\thepage}

\title{Plan: Cross-Correlation Analysis with Optimized Waveform + New Optimization}
\author{BATPING Project}
\date{June 5, 2026}

\begin{document}

\maketitle

\section*{Context}

\textbf{Current state:} We have completed cross-correlation analysis comparing standard xcorr vs GCC-PHAT on \textit{Gaussian waveform} data, showing:
\begin{itemize}
    \item Standard xcorr envelope FWHM: 276--333 $\mu$s depending on preprocessing
    \item GCC-PHAT (raw data): 21 $\mu$s FWHM (94\% improvement, but noisy/spurious peaks)
\end{itemize}

\textbf{Optimized waveform available:} Bayesian optimization achieved 41\% FWHM reduction in the \textit{response waveform} (311 $\mu$s vs 530 $\mu$s baseline). Recordings exist:
\begin{itemize}
    \item \texttt{matched\_filtering\_v2\_june1-2/optimized\_2m\_reflection\_50Hz.wav}
    \item \texttt{matched\_filtering\_v2\_june1-2/optimized\_1.8m\_reflection\_50Hz.wav}
\end{itemize}

\textbf{This plan has two parts:}
\begin{enumerate}
    \item Run same xcorr/GCC-PHAT analysis on optimized waveform recordings
    \item Design new Bayesian optimization targeting xcorr peak sharpness directly
\end{enumerate}

\hrule
\vspace{1em}

\section{Part 1: Cross-Correlation Analysis with Optimized Waveform}

\subsection{Task}
Run the same analysis pipeline we used for Gaussian waveform on the Bayesian-optimized waveform recordings.

\subsection{Data Processing Steps}
\begin{enumerate}
    \item Load optimized wav files (2m template, 1.8m signal)
    \item Create npz files with different preprocessing levels:
    \begin{itemize}
        \item \textbf{Envelope}: Bandpass (35--45kHz) + Hilbert envelope
        \item \textbf{Filtered no-hilbert}: Bandpass only (keeps 40kHz oscillation)
        \item \textbf{Raw}: No preprocessing
    \end{itemize}
    \item Run cross-correlation analysis:
    \begin{itemize}
        \item Standard xcorr: all three preprocessing levels
        \item GCC-PHAT: raw data only (bandpass destroys phase info needed for PHAT)
    \end{itemize}
\end{enumerate}

\subsection{Files to Create}
Working directory: \texttt{matched\_filtering\_v2\_june1-2/}

\textbf{NPZ data files:}
\begin{itemize}
    \item \texttt{optimized\_template\_2m\_envelope.npz}
    \item \texttt{optimized\_signal\_1.8m\_envelope.npz}
    \item \texttt{optimized\_template\_2m\_filterednohilbert.npz}
    \item \texttt{optimized\_signal\_1.8m\_filterednohilbert.npz}
    \item \texttt{optimized\_template\_2m\_raw.npz}
    \item \texttt{optimized\_signal\_1.8m\_raw.npz}
\end{itemize}

\textbf{Output plots:}
\begin{itemize}
    \item \texttt{optimized\_xcorr\_envelope.png} -- Standard xcorr, envelope data
    \item \texttt{optimized\_xcorr\_filterednohilbert.png} -- Standard xcorr, filtered data
    \item \texttt{optimized\_xcorr\_raw.png} -- Standard xcorr, raw data
    \item \texttt{optimized\_gccphat\_raw.png} -- GCC-PHAT, raw data
    \item \texttt{optimized\_gccphat\_standalone.png} -- GCC-PHAT correlation plot alone
\end{itemize}

\subsection{Expected Outcome}
Compare FWHM results:

\begin{table}[h]
\centering
\begin{tabular}{lcc}
\toprule
Preprocessing & Gaussian Waveform & Optimized Waveform \\
\midrule
Envelope (std xcorr) & 276 $\mu$s & ? \\
Filtered (std xcorr) & 219 $\mu$s & ? \\
Raw (std xcorr) & 333 $\mu$s & ? \\
Raw (GCC-PHAT) & 21 $\mu$s & ? \\
\bottomrule
\end{tabular}
\end{table}

\hrule
\vspace{1em}

\section{Part 2: Bayesian Optimization Targeting XCorr Peak Sharpness}

\subsection{Motivation}
The current Bayesian optimizer minimizes \textbf{response waveform FWHM} (piezo output envelope width). But what we actually care about is \textbf{cross-correlation peak sharpness} for TOF measurement. While these are correlated, optimizing the end goal directly may yield better results.

\subsection{Key Insight}
Response waveform FWHM determines \textit{minimum achievable} xcorr peak width, but:
\begin{itemize}
    \item Waveform shape (not just width) affects xcorr sidelobes
    \item Phase coherence across the waveform affects GCC-PHAT performance
    \item Optimizing xcorr directly captures all these factors
\end{itemize}

\subsection{Design Considerations}
\textbf{Objective function options:}
\begin{enumerate}
    \item Standard xcorr envelope FWHM (most robust)
    \item GCC-PHAT peak width (sharper but noisy -- may be unstable for optimization)
    \item Peak-to-sidelobe ratio (PSR) -- measures unambiguity
    \item Weighted combination: \texttt{FWHM + $\lambda$ * (1/PSR)}
\end{enumerate}

\textbf{Measurement protocol:}
\begin{itemize}
    \item Send waveform, record response
    \item Compute xcorr between two consecutive pulses (or template vs new)
    \item Extract FWHM and/or PSR from xcorr result
    \item Return as cost to optimizer
\end{itemize}

\textbf{Trade-offs:}
\begin{itemize}
    \item Slower per iteration (need to compute xcorr, not just envelope FWHM)
    \item May find different optimum than waveform-FWHM optimization
    \item Need stable reference (template) for xcorr -- could use self-correlation
\end{itemize}

\subsection{Implementation Approach}
Modify \texttt{bayesian\_optimizer.py} to:
\begin{enumerate}
    \item Add new cost function: \texttt{xcorr\_fwhm\_cost(waveform)}
    \item Record response, extract segment, compute xcorr with template
    \item Measure envelope FWHM of xcorr peak
    \item Return as optimization target
\end{enumerate}

\subsection{Decision: Standard XCorr Envelope FWHM}
\textbf{Objective function:} Standard cross-correlation envelope FWHM (not GCC-PHAT, not self-correlation)
\begin{itemize}
    \item Most robust and stable for optimization loop
    \item Directly measures what matters for TOF precision
    \item Template-based: correlate new response with fixed template
\end{itemize}

\subsection{Implementation Steps}
\begin{enumerate}
    \item Create initial template by recording baseline Gaussian response
    \item For each optimizer iteration:
    \begin{itemize}
        \item Generate candidate waveform
        \item Send to piezo, record response
        \item Compute xcorr(template, response)
        \item Extract envelope via Hilbert transform
        \item Measure FWHM of envelope peak
        \item Return FWHM as cost (minimize)
    \end{itemize}
\end{enumerate}

\subsection{Files to Modify/Create}
\begin{itemize}
    \item \texttt{bayesian\_optimization/bayesian\_optimizer\_xcorr.py} -- New optimizer targeting xcorr FWHM
    \item Reuse existing hardware interface from \texttt{bayesian\_optimizer.py}
\end{itemize}

\hrule
\vspace{1em}

\section{Verification}

\subsection{Part 1 Verification}
\begin{itemize}
    \item Compare optimized vs Gaussian waveform FWHM numbers
    \item Visual comparison of xcorr peak shapes
    \item Check if optimized waveform gives tighter peaks
\end{itemize}

\subsection{Part 2 Verification}
\begin{itemize}
    \item Run optimization for $\sim$50--100 iterations
    \item Compare best xcorr FWHM vs baseline
    \item Compare with result from waveform-FWHM optimization
    \item Test: does xcorr-optimized waveform give better TOF precision?
\end{itemize}

\hrule
\vspace{1em}

\section{Original Context (Reference)}

BATPING measures dome temperature fluctuations via changes in the speed of sound. The core measurement is time-of-flight (TOF) of an ultrasonic pulse along a fixed path. Temperature changes cause speed-of-sound changes:

$$c(T) = 331.3 \sqrt{1 + T/273.15} \approx 331.3 + 0.6T \text{ m/s}$$

\textbf{Primary Goal:} Verify that the system can achieve \textbf{0.1$^\circ$C RMS temperature resolution}.

\subsection{Experimental Setup (Lab Configuration)}

\textbf{Geometry:} Reflector-based round-trip measurement
\begin{itemize}
    \item Piezo array and microphone positioned side-by-side, both facing a wooden slab reflector
    \item One-way distance to reflector: \textbf{1.406 m}
    \item Round-trip path length: \textbf{2.812 m}
    \item Expected TOF: 2.812 m / 343 m/s $\approx$ \textbf{8.2 ms}
    \item Reflector: flat wooden slab oriented perpendicular to acoustic path
\end{itemize}

\textbf{Why this works:} Both flat reflector (1 reflection) and corner-cube retroreflector (3 reflections) use an odd number of reflections, preserving pressure phase. The flat slab loses the angular tolerance of a CCR but works if aligned orthogonally.

\textbf{Considerations:}
\begin{itemize}
    \item Align slab carefully --- tilt will deflect return beam away from mic
    \item Mic/piezo should be close together to share the same acoustic path
    \item Wood absorbs some energy --- expect lower SNR than CCR, but should be adequate
\end{itemize}

\subsection{Path Length Analysis}

\begin{table}[h]
\centering
\begin{tabular}{lccc}
\toprule
Path length & Shift/1$^\circ$C & Shift/0.1$^\circ$C & Samples @ 384kHz \\
\midrule
18.7 cm (old bench) & 0.95 $\mu$s & 95 ns & 0.037 samples \\
1 m & 5.1 $\mu$s & 510 ns & 0.20 samples \\
\textbf{2.812 m (current)} & \textbf{14.3 $\mu$s} & \textbf{1.43 $\mu$s} & \textbf{0.55 samples} \\
\bottomrule
\end{tabular}
\end{table}

\textbf{Key advantage:} 2.8m round-trip gives 15$\times$ more TOF shift per degree than the 18.7cm bench setup. Achieving 0.1$^\circ$C resolution requires $\sim$1.4 $\mu$s precision --- about 0.55 samples --- achievable with sub-sample interpolation.

\subsection{Verification Approach: Measure TOF Repeatability Floor}

\begin{enumerate}
    \item Fix all conditions (constant temperature)
    \item Take N measurements (100+)
    \item Compute $\sigma_\tau$ = std dev of measured TOF
    \item Convert to equivalent temperature noise:
    $$\sigma_T = \sigma_\tau \cdot \frac{c_0^2}{0.6 \cdot d}$$
    \item Compare $\sigma_T$ to 0.1$^\circ$C target
\end{enumerate}

This characterizes the \textbf{noise floor} of the system. If $\sigma_T > 0.1^\circ$C, options are:
\begin{itemize}
    \item Longer path length
    \item More averaging per measurement
    \item Better correlation method (GCC-PHAT vs standard)
    \item Different input waveform (optimized pulse)
\end{itemize}

\subsection{Approach}

\subsubsection{Input Waveforms to Test}

Compare TOF resolution across different input pulses:

\begin{enumerate}
    \item \textbf{Gaussian $\sigma$=2 cycles} --- narrow, sharp, may have lower SNR
    \item \textbf{Gaussian $\sigma$=5 cycles} --- baseline, broader but higher energy
    \item \textbf{Gaussian $\sigma$=10 cycles} --- wide, high energy but potentially worse resolution
    \item \textbf{Bayesian-optimized pulse} --- from \texttt{bayesian\_optimizer\_state.npz}, achieved 38\% FWHM reduction
\end{enumerate}

Hypothesis: narrower pulses $\rightarrow$ sharper correlation peak $\rightarrow$ better TOF resolution, but may trade off SNR.

\subsubsection{Template Creation}
\begin{enumerate}
    \item Send chosen input pulse to piezo
    \item Record 10 consecutive measurements
    \item Average to create template $h[n]$ --- averaging reduces noise while preserving coherent signal
\end{enumerate}

\subsubsection{TOF Measurement via Matched Filtering}
Cross-correlate template with new measurement:

$$R_{hs}[\tau] = \sum_n h[n] \cdot s[n + \tau]$$

Peak location gives time delay. For sub-sample precision, fit parabola to peak neighborhood.

\subsubsection{GCC-PHAT Option}
Generalized Cross-Correlation with Phase Transform:

$$R_{\text{PHAT}}[\tau] = \mathcal{F}^{-1}\left\{ \frac{H^*(f) S(f)}{|H^*(f) S(f)|} \right\}$$

Whitens spectrum $\rightarrow$ sharper peak, but may amplify noise outside 40kHz band. \textbf{Implement both, compare.}

\subsubsection{GCC-PHAT Peak Ambiguity Problem}

\textbf{Issue:} Pure GCC-PHAT produces impulse-like peaks where the true TOF peak may be only 1--2 samples wide, while noise spikes elsewhere can have comparable amplitude. In the raw data analysis, we observed spurious peaks at $\sim$630 $\mu$s with amplitude comparable to the true peak at $\sim$591 $\mu$s.

\textbf{Risk:} Peak detection could lock onto wrong maximum, causing catastrophic TOF errors (tens of $\mu$s instead of sub-$\mu$s precision).

\textbf{Mitigation strategies to evaluate:}

\begin{enumerate}
    \item \textbf{Two-stage hybrid approach} (recommended)
    \begin{itemize}
        \item Stage 1: Standard xcorr to find coarse peak location (robust, unambiguous)
        \item Stage 2: GCC-PHAT within $\pm$200 $\mu$s window of Stage 1 result (sharp, precise)
        \item Benefits: Gets robustness of xcorr + precision of PHAT
    \end{itemize}

    \item \textbf{Weighted GCC-PHAT (PHAT-$\beta$)}
    \begin{itemize}
        \item Partial whitening: \texttt{weight = |X|$^\beta$} where $0 < \beta < 1$
        \item $\beta \approx 0.3$--$0.5$ trades off some sharpness for reduced noise sensitivity
        \item Tune $\beta$ empirically for best robustness/precision tradeoff
    \end{itemize}

    \item \textbf{Peak-to-sidelobe ratio (PSR) gating}
    \begin{itemize}
        \item Compute ratio: main\_peak / second\_highest\_peak
        \item Reject measurements where PSR $<$ threshold (e.g., $< 1.5$)
        \item Flags ambiguous cases for re-measurement or fallback to xcorr
    \end{itemize}

    \item \textbf{Constrained search} (if prior knowledge available)
    \begin{itemize}
        \item Only search for peak within expected TOF $\pm$ margin
        \item Requires stable setup or tracking of previous measurements
    \end{itemize}
\end{enumerate}

\textbf{Evaluation criteria:}
\begin{itemize}
    \item Repeatability ($\sigma_\tau$) under static conditions
    \item Robustness: \% of measurements with PSR $>$ threshold
    \item Accuracy: mean offset from expected value
\end{itemize}

\subsubsection{Temperature Calculation}
Given TOF shift $\Delta\tau$ relative to template:

$$\Delta T = \frac{c_0^2}{0.6 \cdot d} \cdot \Delta\tau$$

where $d$ is path length, $c_0 \approx 343$ m/s at 20$^\circ$C.

\subsection{Implementation}

\subsubsection{New file: \texttt{agilent\_project/tof\_temperature\_monitor.py}}

\begin{lstlisting}[language=Python]
class TOFTemperatureMonitor:
    def __init__(self, path_length_m, reference_temp_c=20.0):
        self.path_length = path_length_m  # Round-trip distance
        self.T_ref = reference_temp_c
        self.template = None
        self.template_tof = None  # Absolute TOF of template

    def create_template(self, n_avg=10):
        """Record n_avg pulses, average to create template."""

    def measure_tof(self, method='xcorr'):
        """Measure absolute TOF using cross-correlation.
        method: 'xcorr' (standard) or 'gcc-phat'
        Returns: TOF in seconds
        """

    def tof_to_delta_temp(self, tof_s):
        """Convert TOF to temperature change from reference."""

    def run_repeatability_test(self, n_measurements=100):
        """Take N measurements, compute statistics."""
\end{lstlisting}

\subsubsection{Key functions}

\begin{enumerate}
    \item \texttt{gcc\_phat(template, signal, fs)} --- GCC-PHAT cross-correlation
    \item \texttt{xcorr\_peak(template, signal)} --- Standard cross-corr with parabolic interpolation
    \item \texttt{parabolic\_interpolation(y)} --- Sub-sample peak refinement from 3 points
\end{enumerate}

\subsubsection{Reuse from existing code}

\begin{itemize}
    \item \texttt{bayesian\_optimizer.py}: Hardware connection (\texttt{connect}, \texttt{record}, \texttt{extract\_envelope}), bandpass filter
    \item \texttt{agilent\_33220a.py}: \texttt{output\_arbitrary()}, \texttt{generate\_gaussian\_pulse()}
    \item Saved optimized waveform: \texttt{bayesian\_optimizer\_state.npz} $\rightarrow$ \texttt{best\_waveform}
\end{itemize}

\subsubsection{Recording parameters for reflector setup}

\begin{itemize}
    \item \textbf{Recording duration:} 15 ms (captures 8.2 ms TOF + margin)
    \item \textbf{Settling time:} 200 ms (let AWG stabilize)
    \item \textbf{Rep rate:} 50 Hz (20 ms period) --- ensures return pulse (8.2 ms) arrives well before next transmit
    \item \textbf{Waveform points:} 4000 @ 200 kSa/s effective rate (or adjust to achieve 20 ms period)
\end{itemize}

\subsection{Files to create}

\begin{enumerate}
    \item \textbf{Create} \texttt{agilent\_project/tof\_temperature\_monitor.py}
\end{enumerate}

\subsection{Verification}

\begin{enumerate}
    \item \textbf{Signal detection:} Verify reflected pulse is visible above noise floor
    \item \textbf{Template creation:} Record 10 measurements, verify averaging reduces noise
    \item \textbf{Static repeatability:} Create template, take 50--100 measurements, compute $\sigma_\tau$
    \item \textbf{Known displacement (optional):} Move reflector 10cm $\rightarrow$ expect 582 $\mu$s shift (20cm round-trip / 343 m/s)
    \item \textbf{GCC-PHAT vs standard:} Compare peak sharpness and repeatability
    \item \textbf{Waveform comparison:} Test input waveforms ($\sigma$=5 baseline, Bayesian-optimized), report:
    \begin{itemize}
        \item Correlation peak width
        \item Repeatability (std dev of TOF)
        \item SNR of correlation peak
    \end{itemize}
\end{enumerate}

\subsection{Expected Outcomes}

\subsubsection{Primary deliverable: TOF repeatability characterization}

For each configuration (waveform $\times$ correlation method), report:
\begin{itemize}
    \item $\sigma_\tau$ (TOF std dev in $\mu$s)
    \item $\sigma_T$ (equivalent temperature noise in $^\circ$C, given 2m path)
    \item Pass/fail vs 0.1$^\circ$C target
\end{itemize}

\subsubsection{Success criteria}
\begin{itemize}
    \item At least one configuration achieves $\sigma_T < 0.1^\circ$C
    \item With 2.812m path, need $\sigma_\tau < 1.43$ $\mu$s for 0.1$^\circ$C resolution
\end{itemize}

\subsubsection{Expected TOF}
\begin{itemize}
    \item Round-trip 2.812m at 343 m/s $\rightarrow$ \textbf{$\sim$8.2 ms} expected arrival time
    \item Recording duration: 15 ms to capture pulse with margin
\end{itemize}

\subsubsection{Output}
\begin{itemize}
    \item Table comparing waveform/method combinations
    \item Recommendation for final system configuration
    \item Plots: correlation peaks, TOF time series, histogram of TOF measurements
\end{itemize}

\end{document}
